% ══════════════════════════════════════════════════════════════════════════════
%  EXPERIMENT 1 — EXPLORING THE LAB                                    (Lesson 2)
% ══════════════════════════════════════════════════════════════════════════════
\experimentdiv{Exploring the Lab}{Naming lab equipment \textbullet{} Drawing scientific diagrams \textbullet{} Finding the safety equipment}{2}\label{sec:1}
\lessoncodes{VC2S8I02 \textbullet{} VC2S8I03 \textbullet{} VC2S8I08}

\vocabhead
\begin{ctab}
\begin{tabular}{|SW{4.2cm}|W{11.6cm}|}
  \hline\hd & \thead{Word} & \thead{Definition (copy from the slide)} \\ \hline
  & \scaf{\textbf{Hazard}} & \scaf{Anything that could cause harm, such as a flame, hot water or broken glass.} \\ \hline
  & \textbf{Risk} & \\ \hline
  & \textbf{Control} & \\ \hline
  & \textbf{Equipment} & \\ \hline
  & \textbf{Scientific diagram} & \\ \hline
\end{tabular}
\end{ctab}

\begin{lipanel}
\inquiryli{Select the right equipment for a task and draw it as a scientific diagram.}
\understandli{Identify hazards in the lab and the equipment that controls them.}
\sctitle
\begin{itemize}[leftmargin=1.4em, itemsep=0pt, topsep=1pt]
  \item I can name ten pieces of lab equipment and say what each one is used for.
  \item I can draw three pieces of equipment as scientific diagrams.
  \item I can point to the safety equipment in our lab and name the hazard each one controls.
\end{itemize}
\end{lipanel}

\begin{skillsfocus}
\focus{PC7}{Name ten pieces of equipment and choose the right one for a job.}
\focus{PC3}{Find the safety equipment and match each piece to the hazard it controls.}
\focus{CO2}{Show equipment as scientific diagrams: the way scientists draw an experiment.}
\focus{PC1}{Work through the stations in order and follow each instruction card.}
\end{skillsfocus}

% ─── TECHNIQUE ────────────────────────────────────────────────────────────────
\exppart{Technique spotlight}{Blue}

\begin{technique}{Drawing a scientific diagram}{CO2}
A scientific diagram shows equipment as a simple, flat outline, so anyone can see quickly how an experiment is set up. Follow these six rules every time.
\begin{enumerate}[leftmargin=1.6em, itemsep=1pt, topsep=3pt]
  \item Use a \textbf{sharp pencil}, never a pen.
  \item Use a \textbf{ruler} for every straight line.
  \item Draw the equipment \textbf{side-on and flat (2D)}: no 3D shapes, no shading, no colouring in.
  \item Leave the top of an open container \textbf{open}.
  \item Label each part with a \textbf{ruled line} that touches the part. Write labels at the side. Label lines never cross.
  \item Give the diagram a \textbf{title} that says what it shows.
\end{enumerate}
\end{technique}

\begin{center}
\begin{tikzpicture}[scale=0.65]
  % a realistic sketch: 3D rim, shading
  \begin{scope}
    \shade[left color=LightGrey, right color=MidGrey!80] (0,0) arc[start angle=180,end angle=360,x radius=0.9,y radius=0.22] -- (1.8,2.5) arc[start angle=0,end angle=-180,x radius=0.9,y radius=0.22] -- cycle;
    \draw[Ink, line width=0.8pt] (0.9,2.5) ellipse (0.9 and 0.22);
    \draw[Ink, line width=0.8pt] (0,2.5) -- (0,0) arc[start angle=180,end angle=360,x radius=0.9,y radius=0.22] -- (1.8,2.5);
    \fill[Sky!45] (0.02,1.2) arc[start angle=180,end angle=360,x radius=0.88,y radius=0.2] -- (1.78,0.02) arc[start angle=0,end angle=-180,x radius=0.88,y radius=0.2] -- cycle;
    \node[font=\footnotesize\bfseries, text=OrangeDark] at (0.9,-0.75) {\ding{55}\ A sketch, not a scientific diagram};
  \end{scope}
  % the scientific diagram
  \begin{scope}[xshift=7cm]
    \draw[Sky!60, line width=0.6pt] (0,1.2) -- (1.8,1.2);
    \eqBeaker
    \eqlab{(1.8,2.2)}{(2.6,2.2)}{beaker}
    \eqlab{(1.2,1.2)}{(2.6,1.2)}{water}
    \node[font=\footnotesize\bfseries, text=TealGreen] at (1.25,-0.75) {\ding{51}\ A scientific diagram};
  \end{scope}
\end{tikzpicture}
\expfigcap{The same beaker of water drawn two ways.}
\end{center}

\begin{workedex}{A conical flask of water}
\begin{minipage}[c]{0.5\linewidth}
\textbf{Task:} Draw a conical flask half full of water as a scientific diagram.

\smallskip
\textbf{How:} The flask is drawn side-on in ruled pencil lines. The water is a single ruled line across the flask, at its level. The top of the flask is left open. Each label sits at the side on a ruled line that touches the part.
\end{minipage}\hfill
\begin{minipage}[c]{0.44\linewidth}\centering
\begin{tikzpicture}
  \draw[Sky!60, line width=0.6pt] (0.32,1.0) -- (1.58,1.0);
  \eqFlask
  \eqlab{(1.25,2.5)}{(2.4,2.5)}{conical flask}
  \eqlab{(1.2,1.0)}{(2.4,1.0)}{water}
  \node[font=\footnotesize\bfseries] at (1.0,-0.45) {A conical flask of water};
\end{tikzpicture}
\end{minipage}
\end{workedex}

\begin{wedo}{A filter funnel in a conical flask}
Draw a filter funnel sitting in the neck of a conical flask as a scientific diagram. Label both pieces of equipment and give your diagram a title.
\drawbox{4.5cm}{}
\end{wedo}

% ─── PLAN AND SAFETY / METHOD ────────────────────────────────────────────────
\exppart{Plan, safety and method}{Teal}

\textbf{Aim:} To identify common laboratory equipment and the safety equipment in our lab, and to draw equipment as scientific diagrams.\skr{PC1}

\begin{infobox}{Safety}
Do not light, heat or plug in anything during this lesson. Carry glassware with two hands and put it down away from the edge of the bench. Report any chipped or broken glass to your teacher: do not pick it up.
\end{infobox}

\subsubsection*{Method --- Part 1: Equipment stations}
\begin{enumerate}[itemsep=2pt]
  \item Your teacher has set out ten numbered stations around the room, each with one piece of equipment and an instruction card.
  \item Start at the station your teacher gives you. Move to the next number when your teacher says to.
  \item At each station, find the matching drawing in the table in your results, then write the equipment's name and what it is used for.
  \item Handle each piece carefully and leave it exactly as you found it.
\end{enumerate}

\subsubsection*{Method --- Part 2: Safety equipment hunt}
\begin{enumerate}[itemsep=2pt]
  \item With your group, walk around the lab and find each piece of safety equipment listed in your results. Do not touch the safety shower, eyewash or fire equipment.
  \item Record where each one is, and the hazard it controls.
  \item Back at your bench, draw a map of the lab and mark where each piece of safety equipment is.
\end{enumerate}

\begin{scaffold}{Word bank --- equipment}
beaker \textbullet{} conical flask \textbullet{} test tube \textbullet{} measuring cylinder \textbullet{} filter funnel \textbullet{} evaporating basin \textbullet{} tripod and gauze mat \textbullet{} Bunsen burner \textbullet{} thermometer \textbullet{} retort stand and clamp
\end{scaffold}

% ─── RESULTS ─────────────────────────────────────────────────────────────────
\exppart{Results}{Amber}

\subsubsection*{Results --- Part 1: Equipment stations \hfill \sk{PC7}}
\newcommand{\eqcell}[2][0.5]{\raisebox{-0.5\height}{\eqpic[#1]{#2}}}
\begin{ctab}
\begin{tabular}{|SZ{1.0cm}|Z{2.6cm}|W{4.0cm}|W{7.6cm}|}
  \hline\hd & \thead{No.} & \thead{Diagram} & \thead{Name} & \thead{What it is used for} \\ \hline
  & \scaf{e.g.} & \eqcell{\eqBeaker} & \scaf{Beaker} & \scaf{Holding, mixing and heating liquids. Its markings only give a rough volume.} \\ \hline
  & 1 & \eqcell{\eqFlask} & & \\ \hline
  & 2 & \eqcell{\eqTestTube} & & \\ \hline
  & 3 & \eqcell[0.42]{\eqCylinder} & & \\ \hline
  & 4 & \eqcell{\eqFunnel} & & \\ \hline
  & 5 & \eqcell{\eqBasin} & & \\ \hline
  & 6 & \eqcell{\eqTripod} & & \\ \hline
  & 7 & \eqcell{\eqBunsen} & & \\ \hline
  & 8 & \eqcell[0.42]{\eqThermometer} & & \\ \hline
  & 9 & \eqcell[0.42]{\eqStand} & & \\ \hline
\end{tabular}
\end{ctab}

\needspace{12\baselineskip}
\task Choose three pieces of equipment from the stations. Draw each one as a scientific diagram, following all six rules. Label and title each diagram.\skr{CO2}

\noindent
\begin{minipage}[t]{0.32\linewidth}\drawbox{6cm}{}\nameline{Title:}\end{minipage}\hfill
\begin{minipage}[t]{0.32\linewidth}\drawbox{6cm}{}\nameline{Title:}\end{minipage}\hfill
\begin{minipage}[t]{0.32\linewidth}\drawbox{6cm}{}\nameline{Title:}\end{minipage}

\clearpage
\subsubsection*{Results --- Part 2: Safety equipment hunt \hfill \sk{PC3}}
\begin{ctab}
\begin{tabular}{|SW{4.0cm}|W{4.6cm}|W{6.6cm}|}
  \hline\hd & \thead{Safety equipment} & \thead{Where it is in our lab} & \thead{The hazard it controls} \\ \hline
  & \scaf{Fire blanket} & \scaf{On the wall beside the front door} & \scaf{Fire: it smothers a small fire, or clothing that has caught fire} \\ \hline
  & Fire extinguisher & & \\ \hline
  & Eyewash station & & \\ \hline
  & Safety shower & & \\ \hline
  & First-aid kit & & \\ \hline
  & Emergency gas shut-off & & \\ \hline
  & Safety glasses & & \\ \hline
  & Broken-glass bin & & \\ \hline
  & Spill kit & & \\ \hline
\end{tabular}
\end{ctab}

\needspace{12\baselineskip}
\task Draw a bird's-eye map of the lab: the walls, doors and benches. Mark each piece of safety equipment with its first letter, and add a key.
\drawbox{8cm}{}

% ─── ANALYSIS ────────────────────────────────────────────────────────────────
\exppart{Analysis and evaluation}{Coral}

\task Explain the difference between a \emph{hazard} and a \emph{risk}. Use a lit Bunsen burner as your example.
\answerbox{5cm}

\needspace{12\baselineskip}
\task Complete the table: for each hazard, name the safety equipment you would use, or the rule that prevents it.\skr{PC3}
\begin{ctab}
\begin{tabular}{|SW{6.6cm}|W{8.6cm}|}
  \hline\hd & \thead{Hazard} & \thead{Safety equipment or rule that controls it} \\ \hline
  & \scaf{Long hair near a Bunsen flame} & \scaf{Tie long hair back before lighting a Bunsen burner.} \\ \hline
  & A chemical splashes into someone's eye & \\ \hline
  & A beaker is knocked off the bench and breaks & \\ \hline
  & Water is spilt on the floor & \\ \hline
  & A student's sleeve catches fire & \\ \hline
\end{tabular}
\end{ctab}

\needspace{10\baselineskip}
\task Scientists draw equipment as flat outlines instead of realistic pictures. Give two reasons why.\skr{CO2}
\answerbox{5cm}

\needspace{12\baselineskip}
\task Name the best piece of equipment for each job. Choose from: beaker, measuring cylinder, filter funnel, test-tube holder, thermometer, evaporating basin.\skr{PC7}
\pt{a}Measuring exactly 35\,mL of water.
\answerbox{1cm}
\pt{b}Heating 100\,mL of water.
\answerbox{1cm}
\pt{c}Separating sand from water.
\answerbox{1cm}
\pt{d}Holding a test tube above a flame.
\answerbox{1cm}
\pt{e}Measuring how hot some water is.
\answerbox{1cm}

\begin{yr8box}
\begin{keepext}
\exttask A measuring cylinder is never heated over a Bunsen burner, but a beaker often is. Suggest why, thinking about what each piece of equipment is designed to do well.
\answerbox{5cm}
\end{keepext}
\end{yr8box}

% ─── SKILL CHECK ─────────────────────────────────────────────────────────────
\exppart{Skill check}{NavyBlue}

\begin{skillcheck}
\scrow{PC7}{I can name ten pieces of lab equipment and say what each one is used for.}
\scrow{CO2}{I can draw equipment as a scientific diagram: pencil, ruler, 2D outline, ruled labels and a title.}
\scrow{PC3}{I can point to each piece of safety equipment in our lab and name the hazard it controls.}
\scrow{PC1}{I followed the station instructions in order and left each station as I found it.}
\end{skillcheck}

\task Which piece of equipment are you least sure about? Write what you will do to learn it before the next practical.
\answerbox{5cm}
